% Plantilla de informe tipo IEEE (conferencia) en español.
% Para compilar con la clase oficial en Overleaf: cambia la primera línea por \documentclass[conference]{IEEEtran}
% y quita \usepackage{ieee_local} y la línea \babelprovide (usa babel [spanish]).
\documentclass[10pt,twocolumn,letterpaper]{article}
\usepackage[utf8]{inputenc}
\usepackage[T1]{fontenc}
\usepackage{mathptmx}
\usepackage[english]{babel}
\babelprovide[import,main,hyphenrules=nohyphenation]{spanish}
\addto\captionsspanish{\renewcommand{\figurename}{Fig.}\renewcommand{\tablename}{Tabla}\renewcommand{\refname}{Referencias}}
\emergencystretch=3em \tolerance=1500
\usepackage{amsmath}\usepackage{graphicx}\usepackage{booktabs}\usepackage{cite}\usepackage{url}\usepackage{microtype}
\usepackage{ieee_local}
\renewcommand{\IEEEkeywordsname}{Palabras clave}

\begin{document}
\title{T\'itulo del trabajo}
\author{\IEEEauthorblockN{Nombre Apellido \hspace{1.5em} Nombre Apellido}
\IEEEauthorblockA{\textit{Facultad o instituto}\\
\textit{Universidad}\\
Ciudad, Pa\'is}}
\maketitle

\begin{abstract}
Resumen de 150 a 200 palabras: problema, m\'etodo, resultado principal y l\'imite del trabajo.
\end{abstract}
\begin{IEEEkeywords}
palabra uno, palabra dos, palabra tres
\end{IEEEkeywords}

\section{Introducci\'on}
Contexto y por qu\'e importa. Cita as\'i \cite{ejemplo}.

\section{M\'etodo}
\subsection{Datos}
Describe los datos y su fuente.
\begin{table}[!t]
\caption{T\'itulo de la tabla}
\centering\small
\begin{tabular}{@{}lcc@{}}\toprule
Columna A & Columna B & Columna C\\\midrule
1 & 2 & 3\\\bottomrule
\end{tabular}
\end{table}

\section{Resultados}
% \begin{figure}[!t]\centering\includegraphics[width=0.95\columnwidth]{figura.png}\caption{Pie de figura.}\end{figure}

\section{Conclusiones}
Qu\'e se encontr\'o y qu\'e limita el resultado.

\begin{thebibliography}{99}
\bibitem{ejemplo} A. Autor, ``T\'itulo del art\'iculo,'' \textit{Revista}, vol. 1, n.\textsuperscript{o} 1, pp. 1--10, 2026.
\end{thebibliography}
\end{document}
